\emph{Preliminaries.} By Theorem \claimref{tower.ladder-pairs-satisfy-car-vectors} the modes $a_{x}^{*} = v_{x}$, $a_{x} = w_{x}$, $x = (k, \alpha)$, satisfy $\{a_{x}, a_{y}^{*}\} = \delta_{xy}$, $\{a_{x}, a_{y}\} = 0$, $v_{x}^{*} = w_{x}$. Hence $[a_{x}^{*}a_{y}, a_{z}^{*}a_{u}] = \delta_{yz}a_{x}^{*}a_{u} - \delta_{ux}a_{z}^{*}a_{y}$, and for coefficient arrays $A, B$ on the one-particle labels the bilinear elements $\hat A = \sum A_{xy}a_{x}^{*}a_{y}$ satisfy $[\hat A, \hat B] = \widehat{[A, B]}$ and $(\hat A)^{*} = \widehat{A^{*}}$ (coefficients in $\mathrm{span}\{1, J\}$, which commutes with the modes since every plane generator anticommutes with both kernel generators, Theorem \claimref{std-algebra.orthogonal-anticommute}). The density with channel array $M$ and shift $m$ is $\rho_{m}(M) = \widehat{S_{m}\otimes M}$, $(S_{m})_{k,k+m} = 1$ when $k, k + m \in W$ and zero otherwise; $S_{m}^{T} = S_{-m}$, so $\rho_{m}(M)^{*} = \rho_{-m}(M^{*})$.

\emph{(i).} $[S_{m}\otimes M, S_{l}\otimes M'] = S_{m}S_{l}\otimes MM' - S_{l}S_{m}\otimes M'M$, and $(S_{m}S_{l})_{k,k+m+l} = 1$ exactly when $k, k + m, k + m + l \in W$, $(S_{l}S_{m})_{k,k+l+m} = 1$ exactly when $k, k + l, k + l + m \in W$, all other entries zero. This is the general identity. For $m = q$, $l = -q$ and $M'$ replaced by $M'^{*}$, $S_{q}S_{-q}$ and $S_{-q}S_{q}$ are the diagonal projections onto $W \cap (W - q)$ and $W \cap (W + q)$, which gives the equal-shift form; for one channel and $d = 1$ it is Theorem \claimref{hawking.current-commutator-exactly-edge-split}. Translation by $q$ is a bijection of $W \cap (W - q)$ onto $W \cap (W + q)$, so the two sets have the same size; removing their intersection gives $|E_{-}(q)| = |E_{+}(q)|$. For $M = M' = \id$ the common part cancels and $[\rho_{q}, \rho_{q}^{*}] = n(E_{-}(q)) - n(E_{+}(q))$, so for every state $\omega$, $\omega([\rho_{q}, \rho_{q}^{*}]) = \sum_{E_{-}(q)}\bar n_{k} - \sum_{E_{+}(q)}\bar n_{k}$ with $\bar n_{k} = \omega(n_{k})$, the total occupation of momentum $k$.

\emph{(ii).} If $\bar n_{k} = c$ on $E_{-}(q) \cup E_{+}(q)$, the value is $c(|E_{-}(q)| - |E_{+}(q)|) = 0$. For $q \ne q'$, (i) with $m = q$, $l = -q'$ writes $[\rho_{q}, \rho_{q'}^{*}]$ as a sum of bilinears $a_{k,\alpha}^{*}a_{k+q-q',\beta}$ with $k \ne k + q - q'$, whose expectation vanishes in every state diagonal in momentum ($\omega(a_{k,\alpha}^{*}a_{l,\beta}) = 0$ for $k \ne l$), as are all the seas below. A state that fills, at each momentum, one mode $c_{k}^{*} = \sum_{\alpha}u_{k,\alpha}a_{k,\alpha}^{*}$ ($u_{k}$ a unit channel vector) and leaves its orthogonal complement empty has $\bar n_{k} = |u_{k}|^{2} = 1$ at every $k$; this is the lower-band sea of any two-channel $h(k)$ with nondegenerate bands at the strip momenta, in particular the relativistic sea of a two-component field in any dimension, and the $1 + 1$ Dirac sea of $h(k) = k\sigma_{3}$ (channel $1$ filled for $k < 0$, channel $2$ for $k \ge 0$). A sea $F$ with $F \subset W \cap (W - q) \cap (W + q)$ meets neither strip, so $\bar n = 0$ there; this condition depends on $q$ (for $d = 1$, $N = 3$, $F = \{-2, -1, 0\}$ the values at $q = 1, 2, 3$ are $0, 1, 1$). In the uniform cases $\omega([\rho_{q}, \rho_{q}^{*}]) = 0$ at every $q$ and on every window, and for the bounded sea at every $q$ for which it is clear of the strips; with the vanishing forms at $q \ne q'$, the total density modes carry no canonical pair in those states. For the $1 + 1$ Dirac sea one channel alone, $\rho_{q}(E_{11})$, has the value $\min(q, |I| - q)$ on a window $I$ of $|I|$ momenta (centred so that half are negative), and the other channel the opposite value: the two chiralities cancel in the total, while each channel-resolved density alone carries the one-dimensional value.

\emph{(iii).} For $W = \{-N, \dots, N - 1\}^{d}$, $q_{1} \ge 0$, and $F = \{k_{1} < 0\}$ in each of $c$ channels: $W \cap (W \mp q)$ is a product of intervals, and in each transverse direction $j$ both have $2N - |q_{j}|$ points; in the first direction, $\{k_{1} : k_{1}, k_{1} + q_{1} \in I\}$ contains $\min(N, 2N - q_{1})$ negative values and $\{k_{1} : k_{1}, k_{1} - q_{1} \in I\}$ contains $\max(N - q_{1}, 0)$, whose difference is $\min(q_{1}, 2N - q_{1})$. Multiplying gives the value, $0$ when $q_{1} = 0$; $q_{1} < 0$ follows by $\rho_{q}^{*} = \rho_{-q}$. On a vector that agrees with $F$ on $E_{-}(q) \cup E_{+}(q)$ the operator $n(E_{-}) - n(E_{+})$ acts as this number. Along a single null generator ($d = 1$) with the window centred at $0$ this is the edge-clean value of Theorem \claimref{hawking.current-commutator-exactly-edge-split}. A half-space is not needed for a nonzero value: by (i) any sea with unequally occupied strips gives one; counting on $W = \{-3, \dots, 2\}^{2}$ gives $3$ at $q = (1, 0)$ and $5$ at $q = (1, 1)$ for the quadrant $\{k_{1} < 0, k_{2} < 0\}$, and $5$ at $q = (1, 0)$ for $\{k_{1} + k_{2} < 0\}$. What singles out the half-space is that its value depends on $q$ only through $q_{1}$ and the overlaps, and is a scalar on the sea.

\emph{(iv).} With $M = \id$ and $M' = \sigma_{3}$ in (i), $[\rho_{q}, \rho_{q}(\sigma_{3})^{*}] = s_{3}(E_{-}(q)) - s_{3}(E_{+}(q))$ (the common part again cancels), $s_{3}(S) = \sum_{k\in S}a_{k}^{*}\sigma_{3}a_{k}$. In the sea of (ii) the one-particle density at $k$ is the projection onto $u_{k}$, $\omega(a_{k,\alpha}^{*}a_{k,\beta}) = \bar u_{k,\alpha}u_{k,\beta}$, so $\omega(a_{k}^{*}Xa_{k}) = \langle u_{k}, Xu_{k}\rangle$. In the Pauli algebra $\Cl(3,0)$, where the channel arrays $\sigma_{i}$ represent the generators $e_{i}$ and half the matrix trace is the scalar part, the lower band of $h(k) = \sum_{i}k_{i}e_{i}$ is the projector $P_{-} = \tfrac12(1 - \hat k)$, $\hat k = k/|k|$: $Q(\hat k) = 1$ gives $P_{-}^{2} = P_{-}$ and $\hat kP_{-} = -P_{-}$, and $2\langle P_{-}e_{i}\rangle_{0} = -\hat k_{i}$. So $\omega(a_{k}^{*}\sigma_{3}a_{k}) = -k_{3}/|k|$. For $q = (0, \dots, 0, q)$ with $1 \le q < N$, the strip $E_{-}$ is the bottom $q$ layers in the last direction ($k_{3} < 0$) and $E_{+}$ the top $q$ layers ($k_{3} \ge N - q > 0$), $2q(2N)^{d-1}$ points in all, none at $k = 0$, and
\[
\omega([\rho_{q}, \rho_{q}(\sigma_{3})^{*}]) = \sum_{k\in E_{-}\cup E_{+}}\frac{|k_{3}|}{|k|} .
\]
For $P_{+} = \tfrac12(1 + \sigma_{3})$, normal with $P_{+}^{2} = P_{+}$, (i) gives $[\rho_{q}(P_{+}), \rho_{q}(P_{+})^{*}] = \sum_{E_{-}}a_{k}^{*}P_{+}a_{k} - \sum_{E_{+}}a_{k}^{*}P_{+}a_{k}$, with $\omega(a_{k}^{*}P_{+}a_{k}) = \tfrac12(1 - k_{3}/|k|)$; since $|E_{-}| = |E_{+}|$ the constant parts cancel and the value is $S/2$. For $d = 1$ every term of $S$ is $1$: $S = 2q$. For $d = 3$, $|k_{3}| \ge N - q$ on both strips and $|k|^{2} \le 3N^{2} < 4N^{2}$, so each term lies in $((N - q)/(2N), 1]$ and the sum lies in $(4qN(N - q), 8qN^{2}]$. The factor $(2N)^{d-1}$ counts the transverse momenta of the strips: per transverse momentum the value lies in $(q(N - q)/N, 2q]$, of order $q$ as in $d = 1$.

\emph{Checks.} The exact check verifies (i) in the algebra on a $2\times3$ window with one channel and a $2\times2$ window with two channels ($M = \id$, $M' = \id, \sigma_{3}, \sigma_{1}$), the general form with unequal shifts of both signs and $J$-valued arrays on the $2\times2$ two-channel window, and $|E_{-}| = |E_{+}|$ by counting for boxes in $d = 1, 2, 3$ and an irregular window; (ii) in the algebra for the $1 + 1$ Dirac sea and a sea with a rotated channel state on four momenta ($q = 1, 2, 3$), one channel alone giving $\min(q, 4 - q)$, and by counting for a bounded sea and uniform filling in $d = 3$, the $q$-dependence of the bounded-sea hypothesis ($0, 1, 1$), and a vanishing form at $q \ne q'$; (iii) by counting for $N = 2, 3$, $d = 2, 3$, every $q$, and in the algebra on a $4\times2$ window, and the quadrant and tilted counterexamples; (iv) the value $2q$ and the spin-resolved value $q$ in the algebra for $d = 1$ on six momenta ($q = 1, 2$), the Pauli projector at three rational unit vectors, and for $d = 3$, $N = 2, 3, 4$, the strip count and the two termwise bounds as integer inequalities.
